\documentclass[11pt,a4paper]{article}

% --- Kodowanie i język ---
\usepackage[utf8]{inputenc}
\usepackage[T1]{fontenc}
\usepackage[polish]{babel}
\usepackage{lmodern}

% --- Layout ---
\usepackage[a4paper,margin=2.5cm]{geometry}
\usepackage{setspace}
\onehalfspacing
\usepackage{parskip}

% --- Grafika i tabele ---
\usepackage{graphicx}
\usepackage{booktabs}
\usepackage{longtable}
\usepackage{array}
\usepackage{tabularx}

% --- Listingi kodu ---
\usepackage{listings}
\usepackage{xcolor}
\lstset{
  basicstyle=\ttfamily\small,
  keywordstyle=\color{blue!70!black}\bfseries,
  commentstyle=\color{gray},
  stringstyle=\color{green!50!black},
  showstringspaces=false,
  breaklines=true,
  frame=single,
  rulecolor=\color{gray!50},
  numbers=left,
  numberstyle=\tiny\color{gray},
  captionpos=b
}

% --- Hiperłącza i metadane ---
\usepackage{hyperref}
\hypersetup{
  colorlinks=true,
  linkcolor=blue!50!black,
  urlcolor=blue!60!black,
  citecolor=blue!50!black,
  pdftitle={Plan projektu: Inteligentny asystent hodowcy bydła},
  pdfauthor={sulikk}
}

% --- Tytuł ---
\title{
  \vspace{-1.5cm}
  \textbf{Inteligentny asystent hodowcy bydła} \\
  \large Mobilna aplikacja z modułami AI, skanowaniem kolczyków \\ i integracją z systemem IRZplus \\
  \vspace{0.5cm}
  \large \textit{Plan projektu i dokumentacja techniczna}
}
\author{sulikk \\ \small Praca inżynierska}
\date{\today}

\begin{document}

\maketitle

\begin{abstract}
\noindent
Dokument przedstawia plan projektu pracy inżynierskiej polegającej na zaprojektowaniu i implementacji mobilnej aplikacji do zarządzania stadem bydła. Aplikacja redukuje obciążenia administracyjne polskiego hodowcy bydła poprzez skanowanie kolczyków urzędowych (kody kreskowe/QR, OCR, RFID/NFC), automatyzację zgłoszeń do państwowego systemu identyfikacji i rejestracji zwierząt (IRZplus, ARiMR) oraz asystenta opartego o duży model językowy (LLM). Dokument opisuje wymagania funkcjonalne i niefunkcjonalne, architekturę systemu, model danych, stos technologiczny, harmonogram realizacji oraz analizę ryzyk.
\end{abstract}

\tableofcontents
\newpage

% ============================================================
\section{Wprowadzenie}
% ============================================================

\subsection{Kontekst}

Polski hodowca bydła ma obowiązek prowadzenia szczegółowej ewidencji zwierząt i zgłaszania zdarzeń zwierzęcych (urodzenie, przybycie, ubycie, padnięcie) do Agencji Restrukturyzacji i Modernizacji Rolnictwa (ARiMR) w nieprzekraczalnym terminie 7~dni od ich wystąpienia. Zgłoszenia odbywają się poprzez portal IRZplus. W obecnej formie obowiązek ten jest powszechnie postrzegany jako uciążliwy: wymaga ręcznego przepisywania 14-znakowych numerów kolczyków, logowania do osobnego systemu i wypełniania formularzy często w warunkach polowych --- z dala od komputera stacjonarnego.

Równolegle rosną oczekiwania rynkowe: kupujący bydło na rynku wtórnym chce widzieć rzetelne, zweryfikowane dane o zwierzęciu (historia leczeń, waga, BCS, pochodzenie), a hodowca musi te informacje przygotowywać oddzielnie dla każdego kanału sprzedaży.

\subsection{Cel pracy}

Celem pracy jest zaprojektowanie i implementacja mobilnej aplikacji do zarządzania hodowlą bydła, która:

\begin{itemize}
  \item redukuje obciążenia administracyjne dzięki skanowaniu kolczyków i automatyzacji zgłoszeń do IRZplus,
  \item dostarcza pojedyncze źródło prawdy o zwierzęciu --- te same dane służą jednocześnie ewidencji rejestrowej i publikacji ogłoszeń sprzedażowych,
  \item integruje asystenta opartego o LLM z dostępem do narzędzi backendowych i polityką potwierdzania akcji przez użytkownika.
\end{itemize}

\subsection{Zakres pracy}

Praca obejmuje 10~zadań szczegółowych uzgodnionych z promotorem (sekcja \ref{sec:harmonogram}). Zakres MVP \textbf{(must have)}: zadania 1--7 (ewidencja, skanowanie, integracja IRZplus, kolejka offline). Zakres \textbf{should have}: zadanie 8 (asystent LLM). Zakres \textbf{nice to have}: zadanie 9 (moduł giełdowy).

% ============================================================
\section{Wymagania}
% ============================================================

\subsection{Wymagania funkcjonalne}

\begin{enumerate}
  \item[\textbf{F1}] System umożliwia prowadzenie ewidencji zwierząt (dodawanie, edycja, usuwanie, wyszukiwanie).
  \item[\textbf{F2}] System pozwala rejestrować zdarzenia zwierzęce: urodzenie, przybycie, ubycie, padnięcie.
  \item[\textbf{F3}] System pozwala rejestrować leczenia (lek, dawka, lekarz, data, okres karencji).
  \item[\textbf{F4}] Skaner kodów kreskowych/QR rozpoznaje numer kolczyka urzędowego.
  \item[\textbf{F5}] OCR rozpoznaje 14-znakowy numer kolczyka (PL + 12~cyfr) z walidacją cyfry kontrolnej.
  \item[\textbf{F6}] Aplikacja odczytuje elektroniczne kolczyki RFID/NFC zgodne z ISO~11784/11785 przez wbudowany moduł NFC smartfona.
  \item[\textbf{F7}] Klient API IRZplus uwierzytelnia się i pobiera dane stada.
  \item[\textbf{F8}] Klient API IRZplus składa zgłoszenia zdarzeń (urodzenie, przybycie, ubycie, padnięcie).
  \item[\textbf{F9}] Kolejka zgłoszeń obsługuje pracę offline --- zdarzenia trafiają do IRZplus przy pierwszej dostępnej łączności.
  \item[\textbf{F10}] Widok zgodności pokazuje stan realizacji obowiązku 7-dniowego zgłoszenia.
  \item[\textbf{F11}] Asystent LLM prowadzi dialog w języku naturalnym i wywołuje narzędzia backendowe.
  \item[\textbf{F12}] Żadne zgłoszenie do IRZplus ani publikacja ogłoszenia nie wykonują się bez jawnego potwierdzenia użytkownika.
  \item[\textbf{F13}] Moduł publikacji ogłoszeń wstępnie wypełnia formularz danymi z karty zwierzęcia.
  \item[\textbf{F14}] Ogłoszenie posiada odznakę „dane potwierdzone w IRZplus" po weryfikacji wobec rejestru.
\end{enumerate}

\subsection{Wymagania niefunkcjonalne}

\begin{enumerate}
  \item[\textbf{N1}] Aplikacja pracuje offline --- wszystkie operacje ewidencyjne dostępne bez sieci.
  \item[\textbf{N2}] Praca w warunkach polowych: duże elementy dotykowe, wysoki kontrast, obsługa w rękawicach.
  \item[\textbf{N3}] Bezpieczeństwo: tokeny API IRZplus przechowywane w bezpiecznym magazynie (Android Keystore / Secret Manager); klucz API LLM \emph{nigdy} nie trafia do aplikacji mobilnej.
  \item[\textbf{N4}] Czas reakcji UI poniżej 200\,ms dla operacji na lokalnej bazie.
  \item[\textbf{N5}] Synchronizacja przyrostowa --- nie przesyła całej bazy przy każdym połączeniu.
  \item[\textbf{N6}] Dostępność: aplikacja działa od Android 9 (API~28) wzwyż.
  \item[\textbf{N7}] RODO: dane osobowe (hodowca, lekarz weterynarii) szyfrowane lokalnie; możliwość eksportu i usunięcia.
\end{enumerate}

\subsection{Wybrane user stories}

\begin{itemize}
  \item \emph{Jako hodowca chcę zeskanować kolczyk nowo urodzonego cielaka i jednym kliknięciem dodać go do ewidencji, żeby nie przepisywać ręcznie 14~znaków.}
  \item \emph{Jako hodowca chcę, żeby aplikacja automatycznie przygotowała zgłoszenie urodzenia do IRZplus i wymagała ode mnie tylko potwierdzenia.}
  \item \emph{Jako hodowca chcę widzieć listę nadchodzących terminów 7-dniowych, żeby nie przekroczyć żadnego.}
  \item \emph{Jako hodowca chcę powiedzieć asystentowi „zgłoś urodzenie cielaka od PL00533888..." i żeby przeprowadził mnie przez proces.}
  \item \emph{Jako kupujący chcę widzieć w ogłoszeniu odznakę potwierdzającą, że dane zwierzęcia są zgodne z rejestrem państwowym.}
\end{itemize}

% ============================================================
\section{Architektura systemu}
% ============================================================

\subsection{Diagram kontekstowy (C4 poziom 1)}

\begin{center}
\begin{tabular}{|c|}
\hline
\\
\textbf{Hodowca} (Android, wieś, słabe łącze) \\
$\downarrow\uparrow$ \\
\textbf{Aplikacja mobilna} (Kotlin + Compose, SQLite offline-first) \\
$\downarrow\uparrow$ HTTPS \\
\textbf{Backend} (FastAPI, PostgreSQL, S3) \\
$\downarrow$ \hfill $\downarrow$ \hfill $\downarrow$ \\
\textbf{IRZplus} \hspace{0.5cm} \textbf{LLM API} \hspace{0.5cm} \textbf{Moduł giełdowy} \\
(ARiMR) \hspace{0.5cm} (Claude/GPT) \hspace{0.5cm} (wewnętrzny) \\
\\
\hline
\end{tabular}
\end{center}

\subsection{Komponenty aplikacji mobilnej}

\begin{itemize}
  \item \textbf{Warstwa UI:} Jetpack Compose, ekrany podzielone na funkcje: lista zwierząt, karta zwierzęcia, formularze, skaner, czat asystenta, ogłoszenia.
  \item \textbf{Warstwa domeny:} przypadki użycia (UseCase), niezależne od frameworków.
  \item \textbf{Warstwa danych:} repozytoria, lokalna baza Room, klienci HTTP (Retrofit) do backendu i ewentualnie bezpośrednio do IRZplus.
  \item \textbf{Moduł skanera:} CameraX + ML Kit (kody kreskowe, QR, OCR); Android NFC API (tagi ISO~14443/15693).
  \item \textbf{Synchronizacja:} WorkManager, kolejka zdarzeń z polityką retry.
\end{itemize}

\subsection{Komponenty backendu}

\begin{itemize}
  \item \textbf{API REST} (FastAPI): endpointy \texttt{/animals}, \texttt{/events}, \texttt{/treatments}, \texttt{/listings}, \texttt{/sync}.
  \item \textbf{Orkiestrator agenta LLM}: wywołania do dostawcy modelu, deklaracje narzędzi (\texttt{tool\_definitions}), pętla agentowa.
  \item \textbf{Klient IRZplus}: warstwa abstrakcji nad API ARiMR (autoryzacja, zgłoszenia, status).
  \item \textbf{Magazyn obiektowy}: zdjęcia zwierząt (S3-kompatybilny --- MinIO w fazie dyplomowej).
\end{itemize}

\subsection{Model danych}

\begin{longtable}{p{3cm}p{3cm}p{7cm}}
\toprule
\textbf{Encja} & \textbf{Pola kluczowe} & \textbf{Opis} \\
\midrule
\endhead
Animal & id, ear\_tag (PL+12), birth\_date, breed, sex, herd\_id & Karta zwierzęcia, klucz biznesowy = ear\_tag. \\
Event & id, animal\_id, type, date, irz\_status & Zdarzenie zwierzęce; type $\in$ \{birth, arrival, departure, death\}. \\
Treatment & id, animal\_id, drug, dose, vet, date, withdrawal\_until & Wpis o leczeniu z okresem karencji. \\
Measurement & id, animal\_id, type, value, date, photo\_uri & Pomiar (waga, BCS); type $\in$ \{weight, bcs\}. \\
Listing & id, animal\_id, price, photos, irz\_verified & Ogłoszenie sprzedażowe z odznaką weryfikacji. \\
Inquiry & id, listing\_id, buyer, message, created\_at & Zapytanie kupującego o ogłoszenie. \\
SyncQueue & id, entity, op, payload, status, retries & Lokalna kolejka operacji do wysłania (offline-first). \\
\bottomrule
\end{longtable}

\subsection{Strategia offline-first}

Wszystkie operacje zapisu trafiają najpierw do lokalnej bazy SQLite (Room) i jednocześnie do tabeli \texttt{SyncQueue} ze statusem \texttt{pending}. Komponent \texttt{SyncWorker} (WorkManager) cyklicznie próbuje wysłać wpisy do backendu i --- w przypadku zdarzeń wymagających zgłoszenia --- do IRZplus. Każdy wpis ma licznik prób oraz politykę wykładniczego backoffu.

% ============================================================
\section{Stos technologiczny}
% ============================================================

\subsection{Aplikacja mobilna}

\begin{tabularx}{\linewidth}{l X}
\toprule
Język & Kotlin 2.0 \\
UI & Jetpack Compose, Material 3 \\
Baza lokalna & Room (nad SQLite) \\
Kamera/OCR & CameraX, ML Kit (Barcode Scanning, Text Recognition) \\
NFC & Android NFC API \\
Sieć & Retrofit + OkHttp + Kotlinx Serialization \\
DI & Hilt \\
Asynchroniczność & Kotlin Coroutines + Flow \\
Synchronizacja & WorkManager \\
Min SDK & 28 (Android 9) \\
Target SDK & 36 \\
\bottomrule
\end{tabularx}

\subsection{Backend}

\begin{tabularx}{\linewidth}{l X}
\toprule
Język & Python 3.12 \\
Framework & FastAPI + Uvicorn \\
ORM & SQLAlchemy 2.x + Alembic \\
Baza danych & PostgreSQL (produkcja), SQLite (lokalny dev) \\
Walidacja & Pydantic v2 \\
Magazyn obiektowy & MinIO (S3-kompatybilny) \\
Klient HTTP & httpx \\
Klient LLM & oficjalny SDK dostawcy (\texttt{anthropic} lub \texttt{openai}) \\
Testy & pytest, pytest-asyncio \\
\bottomrule
\end{tabularx}

\subsection{Integracja IRZplus}

Klient HTTP oparty na publicznych schematach Swagger udostępnianych przez ARiMR na portalu \texttt{gov.pl}. Uwierzytelnianie zgodne z instrukcją ARiMR. W przypadku braku dostępu do środowiska sandbox --- warstwa abstrakcji z mockiem opartym o schematy Swagger; testy integracyjne na mocku z udokumentowaniem ograniczenia w pracy.

% ============================================================
\section{Harmonogram realizacji}
\label{sec:harmonogram}
% ============================================================

Harmonogram odwzorowuje 10~zadań karty pracy z naciskiem na priorytetyzację MVP (zadania 1--7).

\begin{longtable}{p{0.6cm}p{8.5cm}p{2cm}p{2.5cm}}
\toprule
\textbf{\#} & \textbf{Zadanie} & \textbf{Priorytet} & \textbf{Estymacja} \\
\midrule
\endhead
1 & Analiza prawna, dokumentacja IRZplus, user stories & must & 2~tyg. \\
2 & Projekt architektury i modelu danych & must & 1~tyg. \\
3 & Lokalna baza SQLite + ekrany ewidencji (lista, karta, leczenia, sprzedaże, statystyki) & must & 3~tyg. \\
4 & Skaner kodów kreskowych/QR + OCR numeru kolczyka PL z walidacją cyfry kontrolnej & must & 2~tyg. \\
5 & Odczyt NFC/RFID (ISO~11784/11785) & must & 1~tyg. \\
6 & Klient API IRZplus (auth, dane stada, zgłoszenia zdarzeń) & must & 3~tyg. \\
7 & Kolejka zgłoszeń offline + widok zgodności 7-dniowej & must & 2~tyg. \\
8 & Asystent LLM z narzędziami i polityką potwierdzania akcji & should & 2~tyg. \\
9 & Moduł publikacji ogłoszeń z weryfikacją IRZplus + zapytania kupujących & nice & 2~tyg. \\
10 & Omówienie wyników, walidacja UX, wnioski & must & 2~tyg. \\
\bottomrule
\end{longtable}

% ============================================================
\section{Ryzyka i mitygacje}
% ============================================================

\begin{longtable}{p{5cm}p{9cm}}
\toprule
\textbf{Ryzyko} & \textbf{Mitygacja} \\
\midrule
\endhead
Brak dostępu do sandbox IRZplus & Mock API oparty o publiczne schematy Swagger; warstwa abstrakcji udowodniona testami integracyjnymi na mocku. \\
Brak fizycznych kolczyków RFID LF~134\,kHz & Użycie zwykłych tagów NFC HF jako substytutu; udokumentowanie opcji rozszerzenia o zewnętrzny czytnik Bluetooth zgodny z ISO~11784/11785. \\
Koszty API modelu LLM & Limit zapytań w wersji demonstracyjnej, cache odpowiedzi, jasna informacja o ekonomice rozwiązania. \\
Kwestie RODO i odpowiedzialność przy automatycznym zgłaszaniu & Zasada „agent proponuje --- użytkownik zatwierdza"; pełen audyt log wysłanych zgłoszeń; udokumentowanie jako świadome ograniczenie. \\
Zakres szerszy niż możliwy do wykonania w czasie dyplomu & Priorytetyzacja zadań MVP (1--7); jawna deklaracja zakresu na początku pracy; zadania 8--9 jako rozszerzenia. \\
Niestabilne łącze na terenach wiejskich & Pełna obsługa offline-first; kolejka zdarzeń z retry; tryb manualnej synchronizacji. \\
\bottomrule
\end{longtable}

% ============================================================
\section{Oczekiwane rezultaty}
% ============================================================

\begin{itemize}
  \item Aplikacja mobilna Android (APK) wraz z kodem źródłowym w repozytorium Git.
  \item Backend (FastAPI) z instrukcją uruchomienia, obsługujący synchronizację, klienta IRZplus i orkiestrację agenta LLM.
  \item Dokumentacja techniczna (architektura, API, instrukcja wdrożenia) jako załącznik do pracy.
  \item Praca dyplomowa opisująca problem, przegląd stanu wiedzy, przyjęte rozwiązanie, wyniki testów oraz wnioski i kierunki dalszego rozwoju.
  \item Demo wideo prezentujące przepływy: skanowanie kolczyka $\rightarrow$ dodanie zwierzęcia $\rightarrow$ zgłoszenie do IRZplus $\rightarrow$ publikacja ogłoszenia $\rightarrow$ rozmowa z agentem.
\end{itemize}

\end{document}
